\section{Finite-order waves and real noncancellation}
\label{sec:four-waves}

Fix the metric $h$, the envelope $f$, and the nested sets supplied by
Lemma~\ref{env:gaps}. We construct real quasimodes with arbitrarily
small weighted residual and a quantitative lower bound on their value
and first derivative taken together. This lower bound will allow the
metric correction to divide locally by either the function or its
squared gradient norm. All constants in this section may depend on the
fixed data and on the finite accuracy requested below, but the displayed
powers in the derivative and noncancellation bounds will not.
Differentiation and distances refer to the fixed coordinates on $P$,
unless a metric is specified.

\subsection{The localized packets}

At each $y\in\overline U_3$, choose three linearly independent vectors
$b_j(y)\in a(y)^\perp$, $1\le j\le3$, such that
\begin{equation}
 |b_j(y)|_h=B(y),\qquad
 H_y(a,a)+H_y(b_j,b_j)>0.
 \label{wave:transverse}
\end{equation}
Such triples exist: the strict envelope condition gives a nonempty open
subset of the sphere of radius $B(y)$ in the three-dimensional space
$a(y)^\perp$, and every such open subset contains three linearly
independent vectors. Each triple extends smoothly to a neighborhood,
by projection into $a^\perp$ and normalization. A finite cover of
$\overline U_3$ therefore allows the choices in
\eqref{wave:transverse} to have uniform positive strictness margins and
uniformly positive Gram determinants. We choose one available triple at
each center. No regularity of this choice as a function of the center is
required. Write $b_j^\flat=h_y(b_j,\cdot)$ and
\[
 c_{yj}=\frac{H_y(a,a)+H_y(b_j,b_j)}{|a(y)|_h^2+|b_j(y)|_h^2}.
\]
These numbers have a positive lower bound and a finite upper bound.

\begin{proposition}[Finite-order packets]\label{wave:packets}
For any integers $\ell,J\ge0$, and all sufficiently large integers $k$,
there are packets indexed by
\[
 Y_k=(k^{-1}\mathbb Z^4)\cap U_3,\qquad y\in Y_k,\quad 1\le j\le3,
\]
of the form
\begin{equation}
 p_{yj,k}(x)=e^{k\Phi_{yj}(x)}A_{yj,k}(x)\xi(x-y).
 \label{wave:packet-form}
\end{equation}
Here $\xi$ is a smooth cutoff supported in a fixed small ball, equal to
one in a smaller ball; all translated supporting balls lie in $P$.
The phases are complex polynomials in $x-y$, and the amplitudes are
polynomials of fixed finite degree in $x-y$ and $k^{-1}$. Their
coefficients are uniformly bounded over the centers and indices. They
satisfy
\begin{align}
 \Phi_{yj}(y)&=f(y),&
 d\Phi_{yj}(y)&=df(y)+i b_j^\flat,\notag\\
 \Re\Hess_h\Phi_{yj}(y)&=H_y-c_{yj}h_y,&
 A_{yj,k}(y)&=1.
 \label{wave:packet-jets}
\end{align}
On each supporting ball, for positive constants $c,C$,
\begin{equation}
 -C|x-y|^2\le\Re\Phi_{yj}(x)-f(x)
       \le-c|x-y|^2,
 \qquad
 A_{yj,k}(x)=1+O(|x-y|+k^{-1}).
 \label{wave:gaussian-bounds}
\end{equation}
The amplitude and each of its fixed-order derivatives are uniformly
bounded. For every fixed multi-index $\alpha$,
\begin{equation}
 |\partial^\alpha p_{yj,k}(x)|
 \le C_\alpha k^{|\alpha|}
              e^{kf(x)-ck|x-y|^2},
 \label{wave:packet-derivatives}
\end{equation}
after decreasing $c$ if necessary, and
\begin{equation}
 \sup_{x\in P}e^{-kf(x)}
 \sum_{|\alpha|\le\ell}
 \big|\partial^\alpha(\Delta_h+\Lambda_k)p_{yj,k}(x)\big|
 \le C_{\ell,J}k^{-J}.
 \label{wave:packet-residual}
\end{equation}
The radii, coefficients, constants, and frequency threshold are fixed
before $k$ varies.
\end{proposition}

\begin{proof}
We give a finite Taylor construction using the complex-phase WKB method;
compare \cite[Section~3]{DenckerSjostrandZworski2004}.
Fix a center and an index and
temporarily suppress them. The eikonal equation is
\begin{equation}
 \langle d\Phi,d\Phi\rangle_h=-2,
 \label{wave:eikonal}
\end{equation}
with the complex bilinear extension of the metric pairing. Its constant
term holds because $a\perp b$ and $|b|_h^2=2+|a|_h^2$. The degree-one
condition is
\[
 (\Hess_h\Phi)_y(a+i b,\cdot)=0.
\]
To prescribe the real Hessian, put $U=H_y-c_{yj}h_y$. A real symmetric
form $V$ must then satisfy
\begin{equation}
 V(b,\cdot)=U(a,\cdot),\qquad
 V(a,\cdot)=-U(b,\cdot).
 \label{wave:hessian-compatibility}
\end{equation}
The two rows of a symmetric form have precisely one compatibility
condition, namely $U(a,a)+U(b,b)=0$. Our choice of $c_{yj}$ gives this
identity. Define the remaining block of $V$ on $\{a,b\}^\perp$ to be
zero. Since $a,b$ are orthogonal with lengths bounded above and away
from zero, this constructs $V$ with uniformly bounded coefficients.
Thus $U+iV$ is the required covariant Hessian; the connection converts
it to the ordinary quadratic Taylor coefficients of $\Phi$.

For completeness, the higher homogeneous equations are also explicitly
soluble. In coordinates write $v=a+i b$ and $D_v=v\cdot\partial$.
For a complex homogeneous polynomial $F_q$ of degree $q$, choose a
complex linear function $\zeta$ with $D_v\zeta=1$. For example,
\[
 \zeta(z)=\frac{\sum_i\overline{v_i}z_i}{\sum_i|v_i|^2}.
\]
Then
\begin{equation}
 \mathcal R_qF_q
 =\sum_{r=0}^q\frac{(-1)^r\zeta^{r+1}}{(r+1)!}D_v^rF_q,
 \qquad D_v\mathcal R_qF_q=F_q.
 \label{wave:jet-inverse}
\end{equation}
The identity follows by telescoping after differentiation.
This is a right inverse from degree $q$ to degree $q+1$, with bounded
norm for each fixed degree, uniformly in our data. For $q\ge2$, the
degree-$q$ eikonal equation contains the new degree-$(q+1)$ phase
coefficient only as $2D_v$ applied to it. All other coefficients in
that equation have already been determined. Hence we can match
\eqref{wave:eikonal} through any prescribed finite Taylor degree.

The transport operator for the eigenvalue $\Lambda_k=2k^2+k$ is
\[
 \mathcal T=2\grad_h\Phi\cdot\partial+\Delta_h\Phi+1.
\]
Indeed, without the cutoff,
\begin{equation}
 e^{-k\Phi}(\Delta_h+\Lambda_k)(e^{k\Phi}A_k)
 =k^2E A_k+k\mathcal T A_k+\Delta_h A_k,
 \qquad E=\langle d\Phi,d\Phi\rangle_h+2.
 \label{wave:transport-identity}
\end{equation}
Here is one finite choice of all the truncation orders. Set
\[
 M=J+\ell+1,\qquad d=2J+3\ell+6,\qquad
 n_m=d+2(M-m)\ (0\le m\le M),\qquad N=n_0+1.
\]
Construct $\Phi$ of degree $N$ so that $E$ vanishes to order $N$
at the center. With this phase fixed, successively construct polynomials
$A^{(m)}$ of degree $n_m$, with $A^{(0)}(y)=1$ and
$A^{(m)}(y)=0$ for $m\ge1$, satisfying
\begin{align}
 \mathcal T A^{(0)}&=O(|x-y|^{n_0}),\notag\\
 \mathcal T A^{(m)}+\Delta_h A^{(m-1)}
     &=O(|x-y|^{n_m})\qquad(1\le m\le M).
 \label{wave:finite-transports}
\end{align}
To see that this is possible, solve the homogeneous equations of degrees
$q=0,\ldots,n_m-1$ in order. The new coefficient of $A^{(m)}$ has
degree $q+1$ and enters as $2D_v$ applied to it, so
\eqref{wave:jet-inverse} applies. The source
$\Delta_h A^{(m-1)}$ uses coefficients of the preceding amplitude of
degree at most $q+2\le n_m+1<n_{m-1}$. Those coefficients have already
been retained. The tapering degrees thus ensure that no later
truncation changes an earlier equation. Only finitely many jets of
the smooth coefficients of $h$ are used. This argument is algebraic
and makes no convergence assertion about an infinite formal series.

Set $A_k=\sum_{m=0}^M k^{-m}A^{(m)}$. All its fixed-order derivatives
are bounded uniformly for $k\ge1$, and it has the stated value and
first-order estimate at the center. Every Taylor remainder in
\eqref{wave:finite-transports} vanishes to order at least $d$, and the
eikonal remainder vanishes to still higher order. Uniform bounds for
their differentiated remainders follow from ordinary smooth Taylor
estimates and the finite bounded constructions above.

The function $\Re\Phi-f$ has zero value and differential at the
center. Its ordinary Hessian there is $-c_{yj}h_y$: the connection
terms cancel because its differential is zero. Uniform Taylor bounds
therefore give \eqref{wave:gaussian-bounds} after choosing a sufficiently
small common radius. They also give the expansion, used below,
\begin{equation}
 d(\Re\Phi_{yj}-f)(x)
 =-c_{yj}h_y(x-y,\cdot)+O(|x-y|^2).
 \label{wave:phase-expansion}
\end{equation}
Choose $\xi$ within this radius. Differentiating its product with the
exponential and amplitude proves \eqref{wave:packet-derivatives}.

Finally, for $q\ge0$ the elementary Gaussian bound
\[
 \sup_{r\ge0}r^q e^{-ckr^2}\le C_q k^{-q/2}
\]
converts Taylor vanishing into powers of $k^{-1}$. Through $\ell$
derivatives, the terms in \eqref{wave:transport-identity} that contain
matched errors are bounded in the $e^{kf}$ scale by
\[
 C k^{2+\ell-(d-\ell)/2}\le Ck^{-J-1}.
\]
The last uncancelled term is $k^{-M}\Delta_h A^{(M)}$; its corresponding
bound is $Ck^{\ell-M}=Ck^{-J-1}$. Derivatives of $\xi$ are supported
at a fixed positive distance from the center and contribute an
exponentially small remainder. This proves
\eqref{wave:packet-residual}. Increasing the finite Taylor orders
changes constants and radii, but introduces no new power of $k$ in
\eqref{wave:packet-derivatives}.
\end{proof}

\subsection{A real Gaussian superposition}

Fix a smooth function $\omega$, compactly supported in $U_2$ and equal
to one near $\overline U_1$. For independent standard real Gaussian
variables $G_{yj},G'_{yj}$ define
\begin{equation}
 u=u_k^0+\omega\Re\sum_{y\in Y_k}\sum_{j=1}^3
             (G_{yj}+iG'_{yj})p_{yj,k},
 \qquad W=e^{kf}+e^{kf_0}\quad\hbox{on }P.
 \label{wave:random}
\end{equation}
The added term extends smoothly by zero to $M$. There are $O(k^4)$
real Gaussian coefficients, and their coefficient event
\[
 \mathcal E_k=\{\,|G_{yj}|\le k,\ |G'_{yj}|\le k
                         \text{ for all }y,j\,\}
\]
satisfies $\mathbb P(\mathcal E_k^c)\le Ck^4e^{-k^2/2}$.

\begin{lemma}[Weighted residual bounds]\label{wave:residual}
For any predetermined integers $\ell,J\ge0$, choose the packets of
Proposition~\ref{wave:packets} with derivative order $\ell$ and residual
power $J+6$. On $\mathcal E_k$, the function in
\eqref{wave:random} equals $u_k^0$ outside a compact subset of $U_2$,
its residual $R=(\Delta_h+\Lambda_k)u$ is compactly supported in $U_2$,
and
\begin{align}
 \sup_P W^{-1}|\partial^\alpha u|
      &\le C_\alpha k^{|\alpha|+5}
          &&\text{for every fixed }\alpha,\label{wave:u-derivatives}\\
 \sup_P W^{-1}\sum_{|\alpha|\le\ell}|\partial^\alpha R|
      &\le C_{\ell,J}k^{-J}.
          \label{wave:weighted-residual}
\end{align}
In particular, the exponents in \eqref{wave:u-derivatives} are
independent of $\ell,J$.
\end{lemma}

\begin{proof}
Summing \eqref{wave:packet-derivatives} over $O(k^4)$ packets and using
the coefficient bound $k$ proves \eqref{wave:u-derivatives}, including
the fixed cutoff and the elementary bounds for $u_k^0$. The summed
packet residual in \eqref{wave:weighted-residual} costs at most five
powers of $k$.

There are two other residual terms. The background residual is
\[
 (\Delta_h+\Lambda_k)u_k^0=(\Delta_h-\Delta_{g_0})u_k^0,
\]
supported in $\supp(h-g_0)\subset U_0$. Each fixed-order derivative
is bounded by a fixed power of $k$ times $e^{kf_0}$. The first gap of
Lemma~\ref{env:gaps}, $f>f_0+\sigma$ on $\overline U_0$, makes this
exponentially small relative to $W$. The commutator terms
$[\Delta_h,\omega]p_{yj,k}$ and their derivatives are supported
outside $U_1$. There the opposite gap $f<f_0-\sigma$ makes their
polynomial bounds times $e^{kf}$ exponentially small relative to $W$.
These are all the residual terms. Their supports, and the exact
background equation off $U_0$, prove the support assertion.
\end{proof}

\subsection{Five real directions prevent cancellation}

The correction in the next sections needs a real function whose value
and gradient do not vanish together, with a quantitative lower bound
in the scale $W$. Three transverse phase directions and one displaced
packet supply the five independent real directions needed in dimension
four.

\begin{lemma}[Real noncancellation]\label{wave:noncancellation}
For every fixed choice of the preceding finite accuracy orders, the event
\begin{equation}
 \mathcal G_k=\mathcal E_k\cap
 \left\{\, |(u(x),k^{-1}du(x))|
                 \ge k^{-200}W(x)
                     \text{ for every }x\in\overline U_2\,\right\}
 \label{wave:good-event}
\end{equation}
has probability tending to one as $k\to\infty$. The norm in this
formula is the Euclidean norm on the value and four coordinate
derivatives. The power $200$ is independent of the requested accuracy.
\end{lemma}

\begin{proof}
Put $\mathcal V_k=W^{-1}(u,k^{-1}du)$. First consider a point where
\begin{equation}
 e^{kf(x)}<k^{-8}W(x).
 \label{wave:low-regime}
\end{equation}
On $\mathcal E_k$ the random contribution to the unscaled value and
normalized gradient is at most $Ck^5e^{kf(x)}\le Ck^{-3}W(x)$.
The background gradient bound in Lemma~\ref{env:background-mode} gives
\[
 |(u_k^0,k^{-1}du_k^0)|\ge c e^{kf_0(x)}
                       \ge c(1-k^{-8})W(x).
\]
Thus $|\mathcal V_k(x)|$ is bounded below by a fixed positive number
throughout this regime, for all large $k$.

It remains to estimate small-ball probabilities at a fixed point
$x\in\overline U_2$ with
\begin{equation}
 e^{kf(x)}\ge k^{-8}W(x).
 \label{wave:high-regime}
\end{equation}
The gap $f<f_0-\sigma$ outside $U_1$ implies that, for large $k$,
such a point lies where $\omega=1$ on a neighborhood. Thus no derivative
of $\omega$ enters this calculation. Choose a grid center $y$ within
$Ck^{-1}$ of $x$, and another center $y'$ within $Ck^{-1}$ of
\[
 x+k^{-1/2}\nu,\qquad \nu=\frac{a(x)}{|a(x)|_h},
\]
viewing the tangent vector in the fixed coordinates. The compact
inclusion $\overline U_2\subset U_3$ ensures that both centers belong
to $Y_k$. They are distinct for all large $k$. The cutoff $\xi$ is
one at $x$ for each selected packet.

For the three packets at $y$ and the packet with index $1$ at $y'$,
the positive numbers
\[
 \beta_{zj}=e^{-kf(x)}|p_{zj,k}(x)|\qquad(z=y\text{ or }y')
\]
are bounded above and away from zero. This follows from both sides of
\eqref{wave:gaussian-bounds}, since $|x-z|=O(k^{-1/2})$, and from
$A_{zj,k}(x)=1+O(k^{-1/2})$.

At this fixed evaluation point, rotate the two Gaussian coordinates of
each selected packet by the argument of $p_{zj,k}(x)$. This is a fixed
orthogonal change of its two real coordinates, preserving their
independent standard Gaussian distribution. After factoring out
$\beta_{zj}$, the two columns of the evaluation map in the $e^{kf(x)}$
scale are, up to changing the sign of the second coordinate,
\begin{equation}
 \begin{pmatrix}1\\
 d\Re\Phi_{zj}(x)+k^{-1}\Re(A_{zj,k}^{-1}dA_{zj,k})(x)
 \end{pmatrix},\qquad
 \begin{pmatrix}0\\
 d\Im\Phi_{zj}(x)+k^{-1}\Im(A_{zj,k}^{-1}dA_{zj,k})(x)
 \end{pmatrix}.
 \label{wave:real-columns}
\end{equation}
The rotation is held fixed in computing these derivatives. The
logarithmic amplitude derivatives appearing here are uniformly bounded
at the selected points.

Select the first column for the packet $(y,1)$, the second columns for
$(y,1),(y,2),(y,3)$, and the first column for $(y',1)$. These are five
independent real Gaussian coordinates: the first two come from the
same orthogonally rotated Gaussian pair, and the other three come
from different pairs. The near-center columns have the forms
\[
 (1,df(x)+O(k^{-1})),\qquad
 (0,b_j^\flat(y)+O(k^{-1}))\quad(1\le j\le3).
\]
For the displaced first column, \eqref{wave:phase-expansion} gives
\begin{align*}
 d\Re\Phi_{y'1}(x)-df(x)
 &=-c_{y'1}h_{y'}(x-y',\cdot)+O(k^{-1})\\
 &=k^{-1/2}c_{y'1}h_x(\nu,\cdot)+O(k^{-1}).
\end{align*}
Subtracting the near first column from this column leaves zero in
the value component. Multiplying the resulting column by $k^{1/2}$
leaves the derivative covector $c_{y'1}h_x(\nu,\cdot)+O(k^{-1/2})$.
The other three derivative covectors are uniformly independent in
$a(y)^\perp$. Since $|x-y|=O(k^{-1})$, they remain uniformly
independent together with $h_x(\nu,\cdot)$. This comparison uses
smoothness of $a,h$, and does not compare the triples chosen at
different centers. The determinant of these five amplitude-factored
columns is therefore at least $ck^{-1/2}$ in absolute value.
Restoring the four bounded positive packet amplitudes, with the
amplitude of $(y,1)$ used twice, preserves this lower bound.

Changing the normalization from $e^{kf(x)}$ to $W(x)$ multiplies
each of the five columns by $e^{kf(x)}/W(x)\ge k^{-8}$. Their
determinant in the $W$ scale is consequently at least
$ck^{-40-1/2}$. Condition on all the unused rotated Gaussian
coordinates. The selected five are still independent standard real
Gaussians; the conditioned coordinates and the deterministic background
give an arbitrary translation. Change of variables in their bounded
Gaussian density now yields, unconditionally as well,
\begin{equation}
 \mathbb P\{\,|\mathcal V_k(x)|<r\,\}
       \le C k^{41}r^5\qquad(r>0)
 \label{wave:small-ball}
\end{equation}
at every point in \eqref{wave:high-regime}. This estimate has not
conditioned on the coefficient event $\mathcal E_k$.

Take a net of mesh at most $k^{-208}$ in the compact coordinate set
$\overline U_2$. It has $O(k^{832})$ points. At its high-regime points
apply \eqref{wave:small-ball} with $r=2k^{-200}$. The probability of
any such small value is at most
\[
 C k^{832+41-1000}=Ck^{-127}.
\]
At its low-regime points the preceding deterministic bound applies on
$\mathcal E_k$. Thus, outside an event of probability tending to zero,
all net values have norm at least $2k^{-200}$ and $\mathcal E_k$ holds.

To pass from the net to the whole set, note that
\[
 W^{-1}|\partial^\alpha W|\le C_\alpha k^{|\alpha|},\qquad
 |d\log W|\le Ck.
\]
Together with \eqref{wave:u-derivatives} through order two, these
give $\|d\mathcal V_k\|\le Ck^6$ on $\mathcal E_k$.
The net interpolation error is at most $Ck^{6-208}=Ck^{-202}$,
which is smaller than $k^{-200}$ for all large $k$. This proves
\eqref{wave:good-event}, even when a point and its chosen net point
belong to different regimes. Intersecting the coefficient and net
events used only a union bound, so no assertion about a truncated
Gaussian density is needed.
\end{proof}
